%% GENERATED by tools/gen_latex.py -- edit content/ch18.py instead.
%% Build:  latexmk -pdf ch18.tex     (or: pdflatex ch18.tex, twice)
\documentclass[aspectratio=169,11pt,t]{beamer}
\usepackage{itbpro}

\coursetitle{Best Practices for the Real World}
\coursesubtitle{Past "works okay": automatic hyperparameter search, ensembling, distributed training across devices, and a near-free 2$\times$ from lower-precision arithmetic.}
\coursekicker{Chapter 18}
\coursesource{Chollet \& Watson, Deep Learning with Python 3e --- chapter 18}
\courseduration{3 hours (2 sessions)}
\coursepresenter{Rahman Indra Kesuma, S.Kom., M.Cs.}{Teaching Assistant}
\coursebrand{AI for Professional \textperiodcentered\ ITB \textperiodcentered\ Ch 18}

\title{Best Practices for the Real World}
\author{Rahman Indra Kesuma, S.Kom., M.Cs.}
\date{}

\begin{document}
\covertitle

\framekicker{Learning outcomes}
\begin{frame}[shrink=25]{Session Objectives}
\leadpar{By the end of this session, participants can:}
\begin{isteps}
\item Distinguish \textbf{hyperparameters from parameters}, and say why hyperparameter search cannot use gradient descent.
\item Define a \textbf{search space} with KerasTuner and run a Bayesian optimization search, then retrain the best configurations properly.
\item Explain \textbf{validation-set overfitting} in tuning, and how to avoid trusting a contaminated number.
\item Ensemble models with a \textbf{weighted average}, and explain why diversity matters more than individual quality.
\item Choose between \textbf{data parallelism and model parallelism}, and say what each one is actually for.
\item Use the \textbf{DeviceMesh and LayoutMap} APIs to shard a model across devices under the JAX backend.
\item Explain \textbf{floating-point precision}, and use float16 inference, mixed-precision training, and loss scaling correctly.
\item Apply \textbf{int8 quantization} to a trained model, and explain why the scale-and-unscale trick preserves the result.
\end{isteps}

\vskip 6pt
\vskip 4pt
\reslink{Book}{Chapter 18 --- full text}{https://deeplearningwithpython.io/chapters/chapter18\_best-practices-for-the-real-world}
\reslink{Notebook}{01\_kerastuner\_search.ipynb}{../../notebooks/ch18/01\_kerastuner\_search.ipynb}
\reslink{Notebook}{02\_ensembling\_and\_weights.ipynb}{../../notebooks/ch18/02\_ensembling\_and\_weights.ipynb}
\reslink{Notebook}{03\_data\_and\_model\_parallel\_jax.ipynb}{../../notebooks/ch18/03\_data\_and\_model\_parallel\_jax.ipynb}
\reslink{Notebook}{04\_mixed\_precision\_and\_loss\_scaling.ipynb}{../../notebooks/ch18/04\_mixed\_precision\_and\_loss\_scaling.ipynb}
\reslink{Notebook}{05\_int8\_quantization.ipynb}{../../notebooks/ch18/05\_int8\_quantization.ipynb}
\reslink{Repo}{Author's official notebook}{https://github.com/fchollet/deep-learning-with-python-notebooks/blob/master/chapter18\_best-practices-for-the-real-world.ipynb}
\vskip 2pt
\end{frame}

\coursesection{01}{Getting the most out of your models}{From "works okay" to "wins machine learning competitions".}

\framekicker{Section 18.1.1}
\begin{frame}[shrink=25]{Hyperparameters are the decisions backpropagation cannot make}
{\fontsize{9.4}{12.6}\selectfont\color{fg2}How many layers? How many units or filters in each? \inlinecode{relu} or something else? \inlinecode{BatchNormalization} after this layer? How much dropout?\par}\vskip 4pt
{\fontsize{9.4}{12.6}\selectfont\color{fg2}These architecture-level choices are called \textbf{hyperparameters}, to distinguish them from the model's \textbf{parameters}, which are trained by backpropagation.\par}\vskip 4pt
\begin{iband}[signal]
Experienced engineers build intuition about what works. But \textbf{there are no formal rules}, and initial decisions are almost always suboptimal even with very good intuition. \hilite{It should not be your job to fiddle with hyperparameters all day.}
\end{iband}
\end{frame}

\framekicker{Section 18.1.1}
\begin{frame}[shrink=25]{The optimization loop}
\begin{center}
\includegraphics[width=0.92\linewidth,height=0.60\textheight,keepaspectratio]{figs/ch18-hpo-loop.tex.pdf}
\end{center}
\vskip 3pt{\color{ink3}\fontsize{7.4}{9.4}\selectfont Automatic hyperparameter optimization is an entire field of research, and an important one.\par}
{\fontsize{9.4}{12.6}\selectfont\color{fg2}The key to the whole process is \textbf{step 4} --- the algorithm that analyses the relationship between validation performance and hyperparameter values to choose what to try next. Bayesian optimization, genetic algorithms, plain random search.\par}\vskip 4pt
\end{frame}

\framekicker{Section 18.1.1 $\cdot$ why it is hard}
\begin{frame}[shrink=25]{Three reasons this is not like training weights}
\begin{center}
\includegraphics[width=0.92\linewidth,height=0.60\textheight,keepaspectratio]{figs/ch18-hpo-hard.tex.pdf}
\end{center}
\vskip 3pt{\color{ink3}\fontsize{7.4}{9.4}\selectfont Training weights is easy by comparison: one loss, one backward pass.\par}
{\fontsize{9.4}{12.6}\selectfont\color{fg2}The third point is the one people underestimate. If a run performs 0.2\% better, was that \textbf{a better configuration} or \textbf{a lucky weight initialization}? Without answering that, you are optimizing noise.\par}\vskip 4pt
\end{frame}

\framekicker{Section 18.1.1 $\cdot$ listing 18.1}
\begin{frame}[shrink=25]{KerasTuner: replace a constant with a range}
{\fontsize{9.4}{12.6}\selectfont\color{fg2}Replace hardcoded values like \inlinecode{units=32} with a range of choices. The set of such choices is the \textbf{search space}.\par}\vskip 4pt
\icode{python}{listing 18.1}{listings/ch18-01.py}
{\fontsize{9.4}{12.6}\selectfont\color{fg2}The function takes \inlinecode{hp} and returns a \textbf{compiled model}. Four kinds of range exist: \inlinecode{Int}, \inlinecode{Float}, \inlinecode{Boolean}, \inlinecode{Choice}.\par}\vskip 4pt
\end{frame}

\framekicker{Section 18.1.1 $\cdot$ listing 18.2}
\begin{frame}[shrink=25]{Or subclass HyperModel for something configurable}
{\fontsize{9.4}{12.6}\selectfont\color{fg2}Model constants become constructor arguments.\par}\vskip 4pt
\icode{python}{listing 18.2}{listings/ch18-02.py}
{\fontsize{9.4}{12.6}\selectfont\color{fg2}\inlinecode{build()} is the same function as before. The gain is \textbf{reuse}: one hypermodel, any number of classes.\par}\vskip 4pt
\end{frame}

\framekicker{Section 18.1.1}
\begin{frame}[shrink=25]{The tuner is a for loop with a strategy}
\begin{center}
\includegraphics[width=0.92\linewidth,height=0.60\textheight,keepaspectratio]{figs/ch18-tuner.tex.pdf}
\end{center}
\vskip 3pt{\color{ink3}\fontsize{7.4}{9.4}\selectfont Three built-in tuners: RandomSearch, BayesianOptimization, Hyperband.\par}
{\fontsize{9.4}{12.6}\selectfont\color{fg2}\textbf{BayesianOptimization} makes smart predictions about which new values are likely to perform best, given the outcome of previous choices --- which is what separates it from random search.\par}\vskip 4pt
\end{frame}

\framekicker{Section 18.1.1}
\begin{frame}[shrink=25]{Configuring the tuner}
{\fontsize{9.4}{12.6}\selectfont\color{fg2}Six arguments, four of which decide whether the search is worth running at all.\par}\vskip 4pt
\icode{python}{}{listings/ch18-03.py}
\begin{itable}{>{\raggedright\arraybackslash}p{0.2632\textwidth}>{\raggedright\arraybackslash}p{0.6768\textwidth}}
\thead{Argument} & \thead{What it controls} \\ \midrule
\inlinecode{objective} & The metric to optimize. \textbf{Always a validation metric} --- the goal of the search is models that generalize. \\
\inlinecode{max\_trials} & How many different configurations to try. \\
\inlinecode{executions\_per\_trial} & How many training runs \textbf{per configuration}, averaged --- this is the answer to the noisy-feedback problem. \\
\inlinecode{overwrite} & \inlinecode{True} to start fresh; \textbf{\inlinecode{False} to resume} a crashed search from the trial logs on disk. \\
\end{itable}
\end{frame}

\framekicker{Section 18.1.1 $\cdot$ sidebar}
\begin{frame}[shrink=25]{Which direction is better?}
{\fontsize{9.4}{12.6}\selectfont\color{fg2}For built-in metrics KerasTuner infers the direction --- accuracy should go up, a loss should go down. For a \textbf{custom metric} you must say so yourself.\par}\vskip 4pt
\icode{python}{}{listings/ch18-04.py}
{\fontsize{9.4}{12.6}\selectfont\color{fg2}\inlinecode{name} must match what appears in the epoch logs. Get it wrong and the search runs happily while optimizing nothing.\par}\vskip 4pt
\end{frame}

\framekicker{Section 18.1.1}
\begin{frame}[shrink=25]{Launching the search}
{\fontsize{9.4}{12.6}\selectfont\color{fg2}\inlinecode{search()} takes the same arguments as \inlinecode{fit()} --- it simply passes them down for each new model.\par}\vskip 4pt
\icode{python}{}{listings/ch18-05.py}
\begin{iband}[rose]
\textbf{Never use your test set as validation data here.} You would overfit to it immediately, and your test metrics would stop meaning anything.
\end{iband}
\end{frame}

\framekicker{Section 18.1.1}
\begin{frame}[shrink=25]{A large epoch count plus early stopping}
{\fontsize{9.4}{12.6}\selectfont\color{fg2}\inlinecode{epochs=100} looks reckless until you notice the \inlinecode{EarlyStopping} callback. You \textbf{do not know in advance} how many epochs each configuration will need, so you give a generous budget and let the callback cut each run short.\par}\vskip 4pt
{\fontsize{9.4}{12.6}\selectfont\color{fg2}The MNIST example above runs in a few minutes. With a typical search space and dataset you will find yourself letting a search run \textbf{overnight, or over several days}.\par}\vskip 4pt
\begin{iband}[amber]
If the process crashes, restart it with \inlinecode{overwrite=False} and it resumes from the trial logs on disk. \hilite{Set the directory somewhere durable before you start a multi-day search.}
\end{iband}
\end{frame}

\framekicker{Section 18.1.1 $\cdot$ listing 18.3}
\begin{frame}[shrink=25]{Retraining the winners is a separate job}
{\fontsize{9.4}{12.6}\selectfont\color{fg2}The search gives you configurations, not finished models.\par}\vskip 4pt
\icode{python}{listing 18.3}{listings/ch18-06.py}
{\fontsize{9.4}{12.6}\selectfont\color{fg2}When retraining, \textbf{fold the validation data back into training} --- you are making no more hyperparameter changes, so there is nothing left to evaluate against it. Here that means training on all of MNIST's training split.\par}\vskip 4pt
{\fontsize{9.4}{12.6}\selectfont\color{fg2}But that leaves one parameter unsettled: \textbf{how many epochs}. The aggressive \inlinecode{patience} used during the search saved time; it may also have left models underfitted.\par}\vskip 4pt
\end{frame}

\framekicker{Section 18.1.1}
\begin{frame}[shrink=25]{First, find the best epoch on the validation set}
{\fontsize{9.4}{12.6}\selectfont\color{fg2}One more run per configuration, this time with a \textbf{much higher patience} --- 10 rather than the 5 used during the search.\par}\vskip 4pt
\icode{python}{}{listings/ch18-07.py}
{\fontsize{9.4}{12.6}\selectfont\color{fg2}The high patience is affordable now: this runs a handful of times, not once per trial.\par}\vskip 4pt
\end{frame}

\framekicker{Section 18.1.1}
\begin{frame}[shrink=25]{Then retrain on everything, for 20\% longer}
{\fontsize{9.4}{12.6}\selectfont\color{fg2}With the epoch count settled, train on the full data --- validation folded back in.\par}\vskip 4pt
\icode{python}{}{listings/ch18-08.py}
{\fontsize{9.4}{12.6}\selectfont\color{fg2}The \inlinecode{* 1.2} accounts for the extra data: \textbf{more samples per epoch means the same number of epochs is a longer run}, but also that the model can absorb a little more before it starts overfitting.\par}\vskip 4pt
\end{frame}

\framekicker{Section 18.1.1 $\cdot$ the shortcut, and the warning}
\begin{frame}[shrink=25]{Two things to remember about tuning at scale}
{\fontsize{9.4}{12.6}\selectfont\color{fg2}If you are not worried about slightly underperforming, there is a shortcut --- reload the top models with the weights saved during the search:\par}\vskip 4pt
\icode{python}{}{listings/ch18-09.py}
\begin{iband}[rose]
And the warning: \textbf{validation-set overfitting}. Because you are updating hyperparameters using a signal computed on validation data, you are effectively \hilite{training them on the validation data}, and they will quickly overfit to it. Always keep this in mind.
\end{iband}
\end{frame}
\note{For a professional audience: this is why a genuinely held-out test set, touched once, is a governance control and not just good hygiene.}

\framekicker{Section 18.1.1 $\cdot$ the art of it}
\begin{frame}[shrink=25]{Tuning is automation, not magic}
{\fontsize{9.4}{12.6}\selectfont\color{fg2}Once upon a time people handcrafted the features that went into shallow models. Deep learning automated that --- features are learned from a feedback signal, not hand-tuned. \textbf{In the same way, you should not handcraft architectures.}\par}\vskip 4pt
{\fontsize{9.4}{12.6}\selectfont\color{fg2}But search spaces grow \textbf{combinatorially} with the number of choices, so turning everything into a hyperparameter is far too expensive. You still need to handpick configurations with the potential to yield good metrics.\par}\vskip 4pt
\begin{iband}[signal]
The gain is that your decisions \textbf{graduate}: from micro-decisions (\textit{how many units in this layer?}) to higher-level ones (\textit{should I use residual connections throughout?}). And higher-level decisions \hilite{generalize across tasks and datasets}.
\end{iband}
\end{frame}

\framekicker{Section 18.1.1}
\begin{frame}[shrink=25]{Premade search spaces}
{\fontsize{9.4}{12.6}\selectfont\color{fg2}Following that logic, KerasTuner ships search spaces relevant to broad categories of problems. Pretty much every image classification problem can be solved with the same \textbf{search space template}.\par}\vskip 4pt
\begin{minipage}[t]{0.494\textwidth}
\begin{icardaccent}
\cardh{kt.applications.HyperXception}
A tunable version of the Keras Applications Xception model.
\end{icardaccent}
\end{minipage}%
\hfill
\begin{minipage}[t]{0.494\textwidth}
\begin{icardaccent}
\cardh{kt.applications.HyperResNet}
The same for ResNet. \textbf{Add data, run the search, get a pretty good model.}
\end{icardaccent}
\end{minipage}%
\end{frame}

\framekicker{Section 18.1.2}
\begin{frame}[shrink=25]{Ensembling, and the blind men}
{\fontsize{9.4}{12.6}\selectfont\color{fg2}\textbf{Ensembling} pools the predictions of several models to produce better predictions. Kaggle winners use very large ensembles that inevitably beat any single model, however good.\par}\vskip 4pt
{\fontsize{9.4}{12.6}\selectfont\color{fg2}The assumption: different well-performing models trained independently are likely to be good \textbf{for different reasons}. Each looks at slightly different aspects of the data, getting part of the truth but not all of it.\par}\vskip 4pt
\begin{iquote}
\quotetext{Blind men meeting an elephant: one touches the trunk --- \textit{"it's like a snake"}; another a leg --- \textit{"like a pillar"}. Each has part of the truth. \textbf{Interviewed together, they can tell a fairly accurate story.}}\quotecite{Section 18.1.2}
\end{iquote}
\end{frame}

\framekicker{Section 18.1.2}
\begin{frame}[shrink=25]{Averaging, and then averaging better}
{\fontsize{9.4}{12.6}\selectfont\color{fg2}The easiest pooling is a plain average at inference time.\par}\vskip 4pt
\icode{python}{}{listings/ch18-10.py}
{\fontsize{9.4}{12.6}\selectfont\color{fg2}This only works if the classifiers are \textbf{more or less equally good}. If one is significantly worse, the result may be worse than the best single model.\par}\vskip 4pt
\end{frame}

\framekicker{Section 18.1.2}
\begin{frame}[shrink=25]{Weighted average, weights learned on validation data}
{\fontsize{9.4}{12.6}\selectfont\color{fg2}The smarter pooling weights each model by how good it is.\par}\vskip 4pt
\icode{python}{}{listings/ch18-11.py}
{\fontsize{9.4}{12.6}\selectfont\color{fg2}Better classifiers get a higher weight. To find a good set, use \textbf{random search} or a simple optimization algorithm such as \textbf{Nelder-Mead} --- over the validation data.\par}\vskip 4pt
\begin{iband}[signal]
There are many variants --- averaging an exponential of the predictions, for instance. In general a \textbf{simple weighted average with weights optimized on validation data is a very strong baseline}.
\end{iband}
\end{frame}

\framekicker{Section 18.1.2 $\cdot$ the actual principle}
\begin{frame}[shrink=25]{Diversity is strength}
\begin{center}
\includegraphics[width=0.92\linewidth,height=0.60\textheight,keepaspectratio]{figs/ch18-ensemble.tex.pdf}
\end{center}
\vskip 3pt{\color{ink3}\fontsize{7.4}{9.4}\selectfont If all the blind men touched only the trunk, they would agree that elephants are like snakes --- and stay wrong together.\par}
{\fontsize{9.4}{12.6}\selectfont\color{fg2}In machine learning terms: \textbf{if all your models are biased the same way, the ensemble keeps that bias.} If they are biased differently, the biases cancel and the ensemble is more robust.\par}\vskip 4pt
\end{frame}

\framekicker{Section 18.1.2 $\cdot$ what follows from that}
\begin{frame}[shrink=25]{As good as possible, while as different as possible}
\begin{minipage}[t]{0.494\textwidth}
\begin{icardgood}
\cardh{Worth doing}
Very different \textbf{architectures}, or even different \textbf{brands} of machine learning --- tree-based methods alongside deep networks.
\end{icardgood}
\end{minipage}%
\hfill
\begin{minipage}[t]{0.494\textwidth}
\begin{icardbad}
\cardh{Largely not worth doing}
The same network trained several times from different random initializations. \textbf{Low diversity, tiny improvement.}
\end{icardbad}
\end{minipage}%
{\fontsize{9.4}{12.6}\selectfont\color{fg2}In 2014 Chollet and Andrei Kolev took \textbf{fourth place} in Kaggle's Higgs Boson challenge with an ensemble of tree models and deep networks. One member --- a regularized greedy forest --- had a \textbf{significantly worse score} than the others and was given a small weight.\par}\vskip 4pt
\begin{iband}[signal]
It improved the ensemble \textbf{by a large factor}, precisely because it was so different: it carried information no other model had. \hilite{It is not about how good your best model is; it is about the diversity of your candidates.}
\end{iband}
\end{frame}

\coursesection{02}{Scaling up with multiple devices}{Faster training directly improves the quality of your solutions.}

\framekicker{Section 18.2 $\cdot$ figure 18.1}
\begin{frame}[shrink=25]{The loop of progress, and its next bottleneck}
{\fontsize{9.4}{12.6}\selectfont\color{fg2}From chapter 7: the quality of an idea is a function of how many refinement cycles it has been through. The speed of iteration depends on how fast you can \textbf{set up}, \textbf{run}, and \textbf{analyse} an experiment.\par}\vskip 4pt
\leadpar{As your Keras fluency grows, coding up an experiment stops being the bottleneck. \textbf{The next bottleneck is training speed.}}
{\fontsize{9.4}{12.6}\selectfont\color{fg2}Fast infrastructure means results back in 10 or 15 minutes, and hence dozens of iterations a day. \hilite{Faster training directly improves the quality of your deep learning solutions.}\par}\vskip 4pt
\end{frame}

\framekicker{Section 18.2.1}
\begin{frame}[shrink=25]{Two kinds of parallelism, for two different problems}
\begin{center}
\includegraphics[width=0.92\linewidth,height=0.60\textheight,keepaspectratio]{figs/ch18-parallelism.tex.pdf}
\end{center}
\vskip 3pt{\color{ink3}\fontsize{7.4}{9.4}\selectfont Model parallelism is not a way to speed up ordinary models --- it is a way to train larger ones.\par}
{\fontsize{9.4}{12.6}\selectfont\color{fg2}And you can mix them: split a model across four devices, then replicate that split model twice, for \textbf{eight devices total}.\par}\vskip 4pt
\end{frame}

\framekicker{Section 18.2.1 $\cdot$ data parallelism}
\begin{frame}[shrink=25]{Divide and conquer}
\begin{center}
\includegraphics[width=0.92\linewidth,height=0.60\textheight,keepaspectratio]{figs/ch18-data-parallel.tex.pdf}
\end{center}
\vskip 3pt{\color{ink3}\fontsize{7.4}{9.4}\selectfont Different samples processed in parallel --- hence \textit{data} parallelism.\par}
{\fontsize{9.4}{12.6}\selectfont\color{fg2}In \textbf{inference} you concatenate the sub-batch predictions. In \textbf{training} you average the gradients and update every replica from the average, so all replicas hold identical weights at all times. This is \textbf{synchronous training}; non-synchronous alternatives exist but are less efficient and no longer used.\par}\vskip 4pt
\end{frame}

\framekicker{Section 18.2.1 $\cdot$ the limitation}
\begin{frame}[shrink=25]{Simple, scalable, and it needs the model to fit}
{\fontsize{9.4}{12.6}\selectfont\color{fg2}Data parallelism is a simple and highly scalable way to train faster: get more devices, increase your batch size, and throughput rises accordingly.\par}\vskip 4pt
\begin{iband}[amber]
One limitation, and it is decisive at the frontier: the model must \textbf{fit on a single device}. It is now common to train foundation models with tens of billions of parameters, which \hilite{will not fit on any single GPU}.
\end{iband}
{\fontsize{9.4}{12.6}\selectfont\color{fg2}That is where model parallelism comes in.\par}\vskip 4pt
\end{frame}

\framekicker{Section 18.2.1 $\cdot$ listing 18.4}
\begin{frame}[shrink=25]{A model too large for one device}
{\fontsize{9.4}{12.6}\selectfont\color{fg2}A concrete example to think with: 16,000 features in, 8,000 potentially overlapping categories out, two dense layers.\par}\vskip 4pt
\icode{python}{listing 18.4}{listings/ch18-12.py}
\begin{minipage}[t]{0.494\textwidth}
\istat{\textasciitilde{}1 billion}{parameters in the first Dense}
\end{minipage}%
\hfill
\begin{minipage}[t]{0.494\textwidth}
\istat{\textasciitilde{}512 million}{parameters in the second}
\end{minipage}%
{\fontsize{9.4}{12.6}\selectfont\color{fg2}With two small devices, data parallelism is not available --- the model does not fit on either. What you can do is \textbf{shard} (or \textit{partition}) a single instance across both.\par}\vskip 4pt
\end{frame}

\framekicker{Section 18.2.1 $\cdot$ two ways to shard}
\begin{frame}[shrink=25]{Horizontal and vertical partitioning}
\begin{center}
\includegraphics[width=0.92\linewidth,height=0.60\textheight,keepaspectratio]{figs/ch18-partition.tex.pdf}
\end{center}
\vskip 3pt{\color{ink3}\fontsize{7.4}{9.4}\selectfont Vertical partitioning is almost always the right choice for large models.\par}
\end{frame}

\framekicker{Section 18.2.1 $\cdot$ vertical partitioning by hand}
\begin{frame}[shrink=25]{Splitting one layer across two devices}
{\fontsize{9.4}{12.6}\selectfont\color{fg2}Split the kernel and bias of the first \inlinecode{Dense} in half, compute each half on its own device, and concatenate.\par}\vskip 4pt
\icode{python}{}{listings/ch18-13.py}
{\fontsize{9.4}{12.6}\selectfont\color{fg2}Each device now holds a kernel of shape \textbf{(16000, 32000)} instead of (16000, 64000). Nobody writes this by hand at scale --- but seeing it once makes the \inlinecode{LayoutMap} API that follows read as ordinary code.\par}\vskip 4pt
\end{frame}

\framekicker{Section 18.2.2 $\cdot$ a strong opinion, stated plainly}
\begin{frame}[shrink=25]{Use JAX}
\begin{iquote}
\quotetext{We will only cover the JAX backend, as it is the most performant and most scalable of the various Keras backends, \textbf{by a mile}. If you're doing any kind of large-scale distributed training and you aren't using JAX, you're making a mistake --- and wasting your dollars burning way more compute than you actually need.}\quotecite{Section 18.2.2}
\end{iquote}
{\fontsize{9.4}{12.6}\selectfont\color{fg2}For getting hold of the devices, there are two realistic options: acquire two to eight GPUs and mount them in one machine --- \textit{it will require a beefy power supply} --- or \textbf{rent a multi-GPU VM} on Google Cloud, Azure, or AWS with drivers preinstalled. For anyone not training 24/7, the second.\par}\vskip 4pt
\end{frame}

\framekicker{Section 18.2.2}
\begin{frame}[shrink=25]{Data parallelism is one line}
{\fontsize{9.4}{12.6}\selectfont\color{fg2}Before building your model, add this. That is the whole change.\par}\vskip 4pt
\icode{python}{}{listings/ch18-14.py}
{\fontsize{9.4}{12.6}\selectfont\color{fg2}For finer control, list the devices and pass them explicitly:\par}\vskip 4pt
\icode{python}{}{listings/ch18-15.py}
{\fontsize{9.4}{12.6}\selectfont\color{fg2}Note the ordering requirement: \textbf{before creating the model}. Setting it afterwards silently does nothing.\par}\vskip 4pt
\end{frame}

\framekicker{Section 18.2.2 $\cdot$ what to actually expect}
\begin{frame}[shrink=25]{N GPUs do not give an N$\times$ speedup}
\begin{minipage}[t]{0.3253\textwidth}
\istat{\textasciitilde{}2.0$\times$}{with 2 GPUs}
\end{minipage}%
\hfill
\begin{minipage}[t]{0.3253\textwidth}
\istat{\textasciitilde{}3.8$\times$}{with 4 GPUs}
\end{minipage}%
\hfill
\begin{minipage}[t]{0.3253\textwidth}
\istat{\textasciitilde{}7.3$\times$}{with 8 GPUs}
\end{minipage}%
{\fontsize{9.4}{12.6}\selectfont\color{fg2}Distribution introduces overhead --- in particular, \textbf{merging the weight deltas} from different devices takes time. The efficiency loss grows with the device count.\par}\vskip 4pt
\begin{iband}[amber]
These numbers assume a \textbf{global batch size large enough to keep every GPU at full capacity}. If your batch is too small, the local batch will not keep the GPUs busy and the speedup collapses.
\end{iband}
\end{frame}

\framekicker{Section 18.2.2 $\cdot$ the DeviceMesh API}
\begin{frame}[shrink=25]{A grid of devices, with named axes}
{\fontsize{9.4}{12.6}\selectfont\color{fg2}A \textbf{device mesh} is a grid of devices, organised along axes. Typically one axis handles data parallelism and one handles model parallelism.\par}\vskip 4pt
\icode{python}{}{listings/ch18-16.py}
{\fontsize{9.4}{12.6}\selectfont\color{fg2}Two devices along axis 0, four along axis 1: computation is \textbf{split across four GPUs}, and \textbf{two copies} of that split model each process a different sub-batch. A mesh need not be 2D, but in practice you will only ever see 1D and 2D.\par}\vskip 4pt
\end{frame}

\framekicker{Section 18.2.2 $\cdot$ the LayoutMap API}
\begin{frame}[shrink=25]{Per variable, per dimension: shard or replicate}
{\fontsize{9.4}{12.6}\selectfont\color{fg2}Variables are the frame of reference. Shard or replicate them across the mesh, and the compiler moves the associated computation to the right device.\par}\vskip 4pt
\begin{center}
\includegraphics[width=0.92\linewidth,height=0.60\textheight,keepaspectratio]{figs/ch18-mesh.tex.pdf}
\end{center}
\vskip 3pt{\color{ink3}\fontsize{7.4}{9.4}\selectfont Two independent decisions, made per dimension of each variable.\par}
{\fontsize{9.4}{12.6}\selectfont\color{fg2}\textbf{Replicate} --- every device along that axis sees the same value. \textbf{Shard} --- a (32, 64) variable becomes four chunks of (32, 16), one per device.\par}\vskip 4pt
\end{frame}

\framekicker{Section 18.2.2}
\begin{frame}[shrink=25]{Variable paths, and the rule of thumb}
{\fontsize{9.4}{12.6}\selectfont\color{fg2}A \textbf{variable path} is a string identifier like \inlinecode{"sequential/dense\_1/kernel"} --- a way to refer to a variable without holding the instance. Print them all with:\par}\vskip 4pt
\icode{python}{}{listings/ch18-17.py}
\begin{iband}[signal]
For a simple model, the go-to rule is: \textbf{shard the last dimension along the \inlinecode{"model"} axis, replicate everything else.}
\end{iband}
\end{frame}

\framekicker{Section 18.2.2}
\begin{frame}[shrink=25]{Writing the layout, and turning it on}
{\fontsize{9.4}{12.6}\selectfont\color{fg2}\inlinecode{None} means replicate along this dimension; a mesh axis name means shard across the devices of that axis.\par}\vskip 4pt
\icode{python}{}{listings/ch18-18.py}
{\fontsize{9.4}{12.6}\selectfont\color{fg2}Once the configuration is set, \textbf{no other part of your code changes} --- the model definition and the training code are the same, whether you use \inlinecode{fit()} or your own loop.\par}\vskip 4pt
\end{frame}

\framekicker{Section 18.2.2}
\begin{frame}[shrink=25]{This is the whole of large-scale training}
\leadpar{Assuming the right \inlinecode{LayoutMap}, the few snippets you just saw are \textbf{enough to distribute any large language model training run} --- scaling to as many devices as you have and to arbitrary model sizes.}
{\fontsize{9.4}{12.6}\selectfont\color{fg2}To check what actually happened, inspect the sharding:\par}\vskip 4pt
\ioutput{listings/ch18-19.txt}
{\fontsize{9.4}{12.6}\selectfont\color{fg2}Or visualise it with \inlinecode{jax.debug.visualize\_sharding(value.shape, value.sharding)}. \textbf{Verify the layout before you pay for a long run} --- a silently wrong layout still trains, just slowly.\par}\vskip 4pt
\end{frame}

\framekicker{Section 18.2.2 $\cdot$ sidebar}
\begin{frame}[shrink=25]{tf.data performance under distribution}
\begin{ibullets}
\item Always provide data as a \textbf{\inlinecode{tf.data.Dataset}} for best performance. NumPy arrays work too --- \inlinecode{fit()} converts them to Datasets anyway.
\item Always \textbf{prefetch}: call \inlinecode{dataset.prefetch(buffer\_size)} before passing to \inlinecode{fit()}.
\item If unsure of the buffer size, use \textbf{\inlinecode{dataset.prefetch(tf.data.AUTOTUNE)}} and let it choose.
\end{ibullets}
{\fontsize{9.4}{12.6}\selectfont\color{fg2}The pattern is the same one chapter 8 introduced. It matters far more here: a starved input pipeline turns eight expensive GPUs into eight expensive idle GPUs.\par}\vskip 4pt
\end{frame}

\framekicker{Section 18.2.3}
\begin{frame}[shrink=25]{TPUs: fast enough to be worth the hoops}
{\fontsize{9.4}{12.6}\selectfont\color{fg2}There is a general trend toward hardware designed specifically for deep learning --- \textbf{ASICs}, application-specific integrated circuits. The most prominent is Google's \textbf{Tensor Processing Unit}.\par}\vskip 4pt
\begin{minipage}[t]{0.494\textwidth}
\istat{15$\times$}{TPU v2 against an NVIDIA P100}
\end{minipage}%
\hfill
\begin{minipage}[t]{0.494\textwidth}
\istat{3$\times$}{more cost-effective than GPU, on average}
\end{minipage}%
{\fontsize{9.4}{12.6}\selectfont\color{fg2}TPU v2 is \textbf{free in Colab} (Runtime $\rightarrow$ Change Runtime Type). Google Cloud offers v3 through v5 for serious runs. With the JAX backend, all you need is the same \inlinecode{keras.distribution.set\_distribution(...)} call --- \textbf{before creating your model}.\par}\vskip 4pt
\end{frame}

\framekicker{Section 18.2.3 $\cdot$ two TPU-specific gotchas}
\begin{frame}[shrink=25]{The data pipeline, and step fusing}
\begin{minipage}[t]{0.494\textwidth}
\begin{icardwarn}
\cardh{GCS read speed becomes the bottleneck}
TPUs process batches extremely quickly. If the dataset is small enough, call \textbf{\inlinecode{dataset.cache()}} so it is read from Cloud Storage only once.
\end{icardwarn}
\end{minipage}%
\hfill
\begin{minipage}[t]{0.494\textwidth}
\begin{icardwarn}
\cardh{Small models underutilize the TPU}
Keeping the cores busy can demand batches \textbf{upward of 10,000 samples}. And with enormous batches you must raise the learning rate --- fewer updates, each more accurate.
\end{icardwarn}
\end{minipage}%
{\fontsize{9.4}{12.6}\selectfont\color{fg2}\textbf{Step fusing} is the way out: run several training steps per TPU execution, doing more work between round trips to VM memory.\par}\vskip 4pt
\icode{python}{}{listings/ch18-20.py}
{\fontsize{9.4}{12.6}\selectfont\color{fg2}For a small model that underutilizes the TPU, this alone can be a \textbf{dramatic} speedup --- and it changes nothing else about the training run.\par}\vskip 4pt
\end{frame}

\coursesection{03}{Lower-precision computation}{Up to 2$\times$ faster on almost any model, basically for free.}

\framekicker{Section 18.3.1}
\begin{frame}[shrink=25]{Precision is to numbers what resolution is to images}
{\fontsize{9.4}{12.6}\selectfont\color{fg2}In mathematics the reals form a continuous axis --- there are infinitely many points between any two numbers, and you can always zoom in. \textbf{In computer science this is not true.} There is a finite number of points between 3 and 4.\par}\vskip 4pt
\begin{center}
\includegraphics[width=0.92\linewidth,height=0.60\textheight,keepaspectratio]{figs/ch18-float.tex.pdf}
\end{center}
\vskip 3pt{\color{ink3}\fontsize{7.4}{9.4}\selectfont How many points depends on how many bits you store the number in.\par}
{\fontsize{9.4}{12.6}\selectfont\color{fg2}The practical way to think about it is the \textbf{smallest distance between two numbers you can safely process}.\par}\vskip 4pt
\end{frame}

\framekicker{Section 18.3.1 $\cdot$ sidebar}
\begin{frame}[shrink=25]{Representable numbers are not uniformly spaced}
{\fontsize{9.4}{12.6}\selectfont\color{fg2}A counterintuitive fact: \textbf{larger numbers have lower precision.} There are as many representable values between 2\textasciicircum{}N and 2\textasciicircum{}(N+1) as there are between 1 and 2, for any N.\par}\vskip 4pt
{\fontsize{9.4}{12.6}\selectfont\color{fg2}That falls out of the encoding --- sign, mantissa, exponent:\par}\vskip 4pt
\ioutput{listings/ch18-21.txt}
{\fontsize{9.4}{12.6}\selectfont\color{fg2}So the error incurred converting a number to floating point \textbf{varies with the value}, and tends to grow with its absolute magnitude.\par}\vskip 4pt
\end{frame}

\framekicker{Section 18.3.2}
\begin{frame}[shrink=25]{float16 inference: one line, nearly 2$\times$}
{\fontsize{9.4}{12.6}\selectfont\color{fg2}Every model so far has used single precision --- float32 weights, float32 computation. That is enough to run forward and backward without losing information, particularly for small gradient updates (learning rates around 1e-3, weight updates around 1e-6).\par}\vskip 4pt
{\fontsize{9.4}{12.6}\selectfont\color{fg2}Modern GPUs and TPUs have hardware that runs 16-bit operations \textbf{much faster and with less memory}.\par}\vskip 4pt
\icode{python}{}{listings/ch18-22.py}
{\fontsize{9.4}{12.6}\selectfont\color{fg2}Set it \textbf{before you define your model}. Expect a nearly \textbf{2$\times$ speed boost} on \inlinecode{model.predict()} on GPU and TPU.\par}\vskip 4pt
\end{frame}

\framekicker{Section 18.3.2}
\begin{frame}[shrink=25]{float16 or bfloat16 --- try both}
\begin{itable}{>{\raggedright\arraybackslash}p{0.282\textwidth}>{\raggedright\arraybackslash}p{0.329\textwidth}>{\raggedright\arraybackslash}p{0.329\textwidth}}
\thead{dtype} & \thead{float16} & \thead{bfloat16} \\ \midrule
Exponent bits & 5 & \textbf{8} \\
Mantissa bits & \textbf{10} & 7 \\
Sign bits & 1 & 1 \\
Character & Finer resolution, narrower range & \textbf{Much wider range}, lower resolution \\
\end{itable}
{\fontsize{9.4}{12.6}\selectfont\color{fg2}Table 18.1. \inlinecode{bfloat16} covers a far wider range of values at lower resolution over that range, and \textbf{works better on some devices, particularly TPUs}.\par}\vskip 4pt
\begin{iband}[signal]
Some devices are better optimised for one than the other, so \hilite{try both and settle for whichever turns out fastest}. It is a one-line experiment.
\end{iband}
\end{frame}

\framekicker{Section 18.3.3}
\begin{frame}[shrink=25]{Training in 16 bits does not work --- so use both}
{\fontsize{9.4}{12.6}\selectfont\color{fg2}Setting the default to 16 bits is great for inference. For \textbf{training} there is a significant complication: gradient descent will not run smoothly, because gradient updates around 1e-5 or 1e-6 --- which are common --- are not representable.\par}\vskip 4pt
{\fontsize{9.4}{12.6}\selectfont\color{fg2}\textbf{Mixed-precision training} is the hybrid: 16-bit computation where precision does not matter, 32-bit where numerical stability does --- in particular for gradients and variable updates.\par}\vskip 4pt
\icode{python}{}{listings/ch18-23.py}
{\fontsize{9.4}{12.6}\selectfont\color{fg2}\textbf{Most of the speed benefit of 16-bit, without meaningfully impacting model quality.}\par}\vskip 4pt
\end{frame}

\framekicker{Section 18.3.3 $\cdot$ how it works}
\begin{frame}[shrink=25]{compute\_dtype and variable\_dtype}
\begin{center}
\includegraphics[width=0.92\linewidth,height=0.60\textheight,keepaspectratio]{figs/ch18-mixed.tex.pdf}
\end{center}
\vskip 3pt{\color{ink3}\fontsize{7.4}{9.4}\selectfont Two dtypes per layer. Both default to float32; mixed precision changes only the first.\par}
\begin{iband}[amber]
Some operations are \textbf{numerically unstable in float16} --- notably softmax and crossentropy. To opt a specific layer out, pass \inlinecode{dtype="float32"} to its constructor.
\end{iband}
\end{frame}

\framekicker{Section 18.3.4}
\begin{frame}[shrink=25]{Loss scaling: multiply the loss so gradients survive}
{\fontsize{9.4}{12.6}\selectfont\color{fg2}Under mixed precision the gradients stay in float16, so the limited range can round small gradients \textbf{down to zero} --- and the model stops learning.\par}\vskip 4pt
{\fontsize{9.4}{12.6}\selectfont\color{fg2}Gradient values are proportional to the loss, so the fix is simple: multiply the loss by a large scalar.\par}\vskip 4pt
\icode{python}{}{listings/ch18-24.py}
{\fontsize{9.4}{12.6}\selectfont\color{fg2}\textbf{\inlinecode{LossScaleOptimizer} is usually the better option} --- the right scaling value can change over the course of training.\par}\vskip 4pt
\end{frame}

\framekicker{Section 18.3.5}
\begin{frame}[shrink=25]{Why float8 is not simply the next step down}
{\fontsize{9.4}{12.6}\selectfont\color{fg2}float16 in the forward pass is \textbf{the last level of precision that just works} --- enough bits for all intermediate tensors, with float32 reserved for gradient updates.\par}\vskip 4pt
{\fontsize{9.4}{12.6}\selectfont\color{fg2}At float8 that stops being true; you lose too much information. float8 is still usable in \textit{some} computations, but it requires considerable modifications to the forward pass. \textbf{You cannot simply set \inlinecode{compute\_dtype} to float8 and run.}\par}\vskip 4pt
{\fontsize{9.4}{12.6}\selectfont\color{fg2}Keras has a built-in implementation targeting Transformers, covering only \inlinecode{Dense}, \inlinecode{EinsumDense}, and \inlinecode{Embedding}. It tracks past activation values to rescale activations each step so as to use the full float8 range, and overrides part of the backward pass to do the same for gradients.\par}\vskip 4pt
\end{frame}

\framekicker{Section 18.3.5 $\cdot$ when it pays}
\begin{frame}[shrink=25]{float8 can make your model slower}
\begin{iband}[rose]
That added machinery has a \textbf{computational cost}. If your model is too small or your GPU not powerful enough, the cost exceeds the benefit and you get a \hilite{slowdown, not a speedup}.
\end{iband}
\begin{minipage}[t]{0.494\textwidth}
\begin{icardwarn}
\cardh{Model size}
Viable only for \textbf{very large} models, typically over \textbf{5 billion parameters}.
\end{icardwarn}
\end{minipage}%
\hfill
\begin{minipage}[t]{0.494\textwidth}
\begin{icardwarn}
\cardh{Hardware}
Large, recent GPUs such as the \textbf{NVIDIA H100}.
\end{icardwarn}
\end{minipage}%
{\fontsize{9.4}{12.6}\selectfont\color{fg2}\textbf{float8 is rarely used in practice}, except in foundation model training runs. Worth knowing exists; not worth reaching for.\par}\vskip 4pt
\end{frame}

\framekicker{Section 18.3.6}
\begin{frame}[shrink=25]{int8 quantization: a different trick}
{\fontsize{9.4}{12.6}\selectfont\color{fg2}Take an already-trained model with float32 weights and convert them to a lower-precision dtype --- typically \textbf{int8} --- while preserving the numerical correctness of the forward pass as far as possible.\par}\vskip 4pt
\begin{center}
\includegraphics[width=0.92\linewidth,height=0.60\textheight,keepaspectratio]{figs/ch18-quant.tex.pdf}
\end{center}
\vskip 3pt{\color{ink3}\fontsize{7.4}{9.4}\selectfont matmul is linear, so the final unscaling cancels the initial scaling exactly.\par}
{\fontsize{9.4}{12.6}\selectfont\color{fg2}Any loss of accuracy comes \textbf{only from the rounding} that happens when casting to int8 --- not from the matmul itself.\par}\vskip 4pt
\end{frame}

\framekicker{Section 18.3.6 $\cdot$ abs-max scaling}
\begin{frame}[shrink=25]{Doing it by hand, once}
{\fontsize{9.4}{12.6}\selectfont\color{fg2}Casting naively would be destructive: \inlinecode{[[0.1, 0.9], [1.2, -0.8]]} becomes \inlinecode{[[0, 0], [1, 0]]}. So we spread each tensor across the full [-127, 127] range first.\par}\vskip 4pt
\icode{python}{}{listings/ch18-25.py}
{\fontsize{9.4}{12.6}\selectfont\color{fg2}\inlinecode{+ 1e-7} avoids dividing by zero. Rounding and clipping \textbf{before} casting is more accurate than casting directly.\par}\vskip 4pt
\end{frame}

\framekicker{Section 18.3.6}
\begin{frame}[shrink=25]{How accurate is it?}
{\fontsize{9.4}{12.6}\selectfont\color{fg2}Do the matmul in int8, then cast back and divide by the product of the two scales.\par}\vskip 4pt
\icode{python}{}{listings/ch18-26.py}
\ioutput{listings/ch18-27.txt}
{\fontsize{9.4}{12.6}\selectfont\color{fg2}Accurate to about three decimal places. For a large matmul this saves a great deal of compute --- int8 can be considerably faster than even float16 --- and all we added were \textbf{fast elementwise ops}: abs, max, clip, cast, divide, multiply.\par}\vskip 4pt
\end{frame}

\framekicker{Section 18.3.6}
\begin{frame}[shrink=25]{In practice, one method call}
{\fontsize{9.4}{12.6}\selectfont\color{fg2}Nobody should implement quantization by hand. Like float8, int8 is built into \inlinecode{Dense}, \inlinecode{EinsumDense}, and \inlinecode{Embedding} --- which unlocks int8 inference for \textbf{any Transformer-based model}.\par}\vskip 4pt
\icode{python}{}{listings/ch18-28.py}
{\fontsize{9.4}{12.6}\selectfont\color{fg2}\inlinecode{predict()} and \inlinecode{call()} now run partially in int8. \textbf{The weights are converted in place}, so this is a one-way operation on that model object --- quantize a copy if you still need the float32 version.\par}\vskip 4pt
\end{frame}

\framekicker{Choosing between them}
\begin{frame}[shrink=25]{Which lever to pull, and when}
\begin{itable}{>{\raggedright\arraybackslash}p{0.2444\textwidth}>{\raggedright\arraybackslash}p{0.3572\textwidth}>{\raggedright\arraybackslash}p{0.3384\textwidth}}
\thead{Technique} & \thead{What it buys} & \thead{What it costs} \\ \midrule
\textbf{Hyperparameter search} & The last few points of accuracy & Hours to days of compute; validation-set overfitting risk \\
\textbf{Ensembling} & More than any single model can reach & N$\times$ inference cost; needs genuinely diverse members \\
\textbf{Data parallelism} & Up to \textasciitilde{}7.3$\times$ on 8 GPUs & Model must fit on one device; large batches needed \\
\textbf{Model parallelism} & Models that fit nowhere else & A LayoutMap to design and verify \\
\textbf{Mixed precision} & Near 2$\times$ on training, nearly free & Loss scaling; unstable ops opted out by hand \\
\textbf{int8 quantization} & Faster inference than float16 & A small, measurable accuracy loss \\
\end{itable}
\end{frame}

\framekicker{Common failure modes}
\begin{frame}[shrink=25]{Four ways this chapter goes wrong in practice}
\begin{minipage}[t]{0.494\textwidth}
\begin{icardbad}
\cardh{Tuning against the test set}
Passing test data as \inlinecode{validation\_data} to \inlinecode{search()}. Every number you report afterwards is contaminated, and nothing warns you.
\end{icardbad}
\end{minipage}%
\hfill
\begin{minipage}[t]{0.494\textwidth}
\begin{icardbad}
\cardh{Ensembling identical models}
Same architecture, different seeds. Low diversity, negligible gain, N$\times$ the inference bill.
\end{icardbad}
\end{minipage}%


\vskip 4pt

\begin{minipage}[t]{0.494\textwidth}
\begin{icardwarn}
\cardh{set\_distribution after model creation}
It silently does nothing. Both \inlinecode{set\_distribution} and \inlinecode{set\_dtype\_policy} must come \textbf{before} the model exists.
\end{icardwarn}
\end{minipage}%
\hfill
\begin{minipage}[t]{0.494\textwidth}
\begin{icardwarn}
\cardh{Mixed precision without loss scaling}
Small gradients round to zero in float16 and the model quietly stops learning. Use \inlinecode{LossScaleOptimizer}.
\end{icardwarn}
\end{minipage}%
\end{frame}

\framekicker{Summary}
\begin{frame}[shrink=25]{What has to survive this chapter}
\begin{isteps}
\item \textbf{Hyperparameters cannot be learned by gradient descent} --- the space is discrete, the feedback expensive and noisy.
\item \textbf{KerasTuner replaces constants with ranges}; a tuner is a loop with a strategy, and BayesianOptimization is the strategy worth defaulting to.
\item \textbf{Tuning is automation, not magic.} You still design the search space, and your decisions graduate from micro to architectural.
\item \textbf{Validation-set overfitting is real}: tuning trains hyperparameters on the validation data.
\end{isteps}
\end{frame}

\framekicker{Summary}
\begin{frame}[shrink=25]{What has to survive this chapter (2 of 2)}
\begin{isteps}
\item \textbf{Ensembling works on diversity}, not on individual quality --- a worse model of a different kind can lift the whole ensemble.
\item \textbf{Data parallelism is for speed, model parallelism is for size.} JAX, DeviceMesh, LayoutMap --- and nothing else in your code changes.
\item \textbf{N GPUs do not give N$\times$ speedup}: about 7.3$\times$ on eight, and only if the global batch keeps every device busy.
\item \textbf{Mixed precision is close to free}: float16 compute, float32 variables, loss scaling to keep small gradients alive.
\item \textbf{float8 is not the next step down} --- it needs a rewritten forward pass, 5B+ parameters, and recent hardware to pay for itself.
\item \textbf{int8 quantization} scales into [-127, 127], multiplies, and unscales --- and \inlinecode{model.quantize("int8")} does it for you.
\end{isteps}
\vskip 4pt
\reslink{NOTEBOOK}{04\_mixed\_precision\_and\_loss\_scaling.ipynb}{../../course-slides/notebooks/ch18/04\_mixed\_precision\_and\_loss\_scaling.ipynb}
\reslink{NEXT}{Chapter 19 --- The future of AI}{../ch19/index.html}
\vskip 2pt
\end{frame}

\end{document}
